\par\noindent\textbf{(1)} 摆球坐标为 \(x=l\sin\theta,\ y=a\cos(\Omega t)-l\cos\theta\)。对惯性系Lagrangian \(L=m(\dot x^2+\dot y^2)/2-mgy\) 用Euler方程，或在支点加速度框架下求切向力，得 \(\ddot\theta+[g-a\Omega^2\cos(\Omega t)]\sin\theta/l=0\)。支点加速度为 \(-a\Omega^2\cos(\Omega t)\)。
\par\noindent\textbf{(2)} 取 \(\theta=\Theta(t)+\xi(t)\) 且慢变量 \(\Theta\) 在一周期近似常数。高频部分方程为 \(\ddot\xi=(a\Omega^2/l)\cos(\Omega t)\sin\Theta\)，可取 \(\xi=-(a/l)\sin\Theta\cos(\Omega t)\)。微振幅相对于慢角为 \(O(a/l)\)，其符号由对时间二阶导数检验。
\par\noindent\textbf{(3)} 将 \(\theta=\Theta+\xi\) 代回并对快周期平均，使用 \(\langle\cos^2\Omega t\rangle=1/2\)，得 \(\ddot\Theta+(g/l)\sin\Theta+(a^2\Omega^2/(2l^2))\sin\Theta\cos\Theta=0\)。因此 \(V_{\mathrm{eff}}(\Theta)=mgl(1-\cos\Theta)+(ma^2\Omega^2/4)\sin^2\Theta\)，在慢近似下为守恒有效势。
\par\noindent\textbf{(4)} 倒立点 \(\Theta=\pi\) 的有效势曲率 \(V_{\mathrm{eff}}''(\pi)=-mgl+ma^2\Omega^2/2\)，稳定要求 \(a^2\Omega^2>2gl\)。若频率不够快、振幅不小或接近参数共振，截断平均不受保证；不能据该不等式宣称所有有限频率驱动都稳定。
